% Table 2: Scale-restoration ablation, raw vs. restored retrieval, on two datasets and
% two metrics (RMSSE on M5, MSE on ETTm2 -- metrics are not comparable across rows, only
% within each dataset's own raw-vs-restored pair).
% Source (M5): docs/ablation-report.md, "normalization without restoration" vs.
%              "mean scaling (+restore)", 1,000-series M5 validation split.
% Source (ETTm2): docs/scalerag-native-dev-results.json, test_mechanisms.A_restoration
%                 (frozen config scale=mean, k=20; test split, n=80,199 windows,
%                 paired-bootstrap 95% CI, window-level).
\begin{table}[t]
  \centering
  \caption{Scale restoration is decisive on both datasets. Raw retrieval normalises context
    windows for matching but does not rescale the retrieved continuation; restored retrieval
    additionally rescales candidate$\to$query before use (Sec.~\ref{sec:method}). M5 reports
    RMSSE (lower is better, no CI available for this legacy ablation); ETTm2 reports MSE with
    a paired-bootstrap 95\% CI (2,000 resamples, $n=80{,}199$ windows, 7/7 channels improved).
    Relative improvement $= (\mathrm{raw} - \mathrm{restored}) / \mathrm{raw}$; positive throughout.}
  \label{tab:ablation}
  \begin{tabular}{@{}llccc@{}}
    \toprule
    Dataset & Metric & Raw retrieval & Restored retrieval & Relative improvement \\
    \midrule
    M5 (1k, val)    & RMSSE & 2.7884 & 0.7425 & $+73.4\%$ (3.8$\times$ lower) \\
    ETTm2 (test)    & MSE   & 2.9989 & 0.4376 & $+85.4\%$ \; [$+85.2\%$, $+85.6\%$] \\
    \bottomrule
  \end{tabular}
\end{table}
